\documentclass[12pt,a4paper]{article}
\usepackage[utf8]{inputenc}
\usepackage[T1]{fontenc}
\usepackage{amsmath,amssymb,amsthm}
\usepackage{physics}
\usepackage{geometry}
\usepackage{enumitem}
\usepackage{hyperref}
\usepackage{fancyhdr}
\usepackage{titlesec}

\geometry{margin=1in}
\pagestyle{fancy}
\fancyhf{}
\rhead{Lecture 7}
\lhead{Quantum Mechanics --- The Theoretical Minimum}
\cfoot{\thepage}

\titleformat{\section}{\large\bfseries}{}{0em}{}
\titleformat{\subsection}{\normalsize\bfseries}{}{0em}{}

\newtheorem{definition}{Definition}
\newtheorem{theorem}{Theorem}

\title{\textbf{Lecture 7: Entanglement and the Density Matrix} \\ \large Wave Functions, Reduced Descriptions, and Quantum Correlations}
\author{Based on Leonard Susskind's \textit{The Theoretical Minimum}}
\date{}

\begin{document}
\maketitle
\thispagestyle{fancy}

% ============================================================
\section{1. The Wave Function}
% ============================================================

\begin{definition}[Wave function]
Given a state $\ket{\Psi}$ and a complete orthonormal basis $\{\ket{a,b,c,\dots}\}$ labeled by eigenvalues of commuting observables, the \textbf{wave function} is
\[
    \psi(a,b,c,\dots) \equiv \braket{a,b,c,\dots}{\Psi}.
\]
It is the set of expansion coefficients---or equivalently, the inner products of the state vector with each basis vector.  It encodes the same information as $\ket{\Psi}$; every observable's expectation value can be computed from it.
\end{definition}

\noindent The term ``wave function'' is historical---it grew out of Schr\"odinger's description of particles, where $\psi(x)$ literally describes a wave.  In general, the concept applies to any quantum system and need not involve waves at all.

The state can be reconstructed from its wave function via
\[
    \ket{\Psi} = \sum_{a,b,c,\dots} \psi(a,b,c,\dots)\,\ket{a,b,c,\dots}.
\]

\subsection{1.1 Examples}
\begin{itemize}[nosep]
    \item \textbf{Single spin}: basis $\{\ket{\uparrow},\ket{\downarrow}\}$; wave function $\psi(\sigma_z)$ takes two values.
    \item \textbf{Two spins}: basis $\{\ket{\uparrow\uparrow},\ket{\uparrow\downarrow},\ket{\downarrow\uparrow},\ket{\downarrow\downarrow}\}$; wave function $\psi(\sigma_z,\tau_z)$ is a function of two discrete variables.
    \item \textbf{Particle on a line}: basis $\{\ket{x}\}$; wave function $\psi(x)$ is a continuous complex-valued function.
\end{itemize}

% ============================================================
\section{2. Expectation Values via the Wave Function}
% ============================================================

For an observable $L$ with matrix elements $L_{a'a}=\bra{a'}L\ket{a}$ (restricting to a single label for clarity):
\[
    \expval{L} = \sum_{a,a'} \psi^*(a')\,L_{a'a}\,\psi(a).
\]
This is the master formula from which all measurable predictions follow.

% ============================================================
\section{3. Composite Systems and the Wave Function}
% ============================================================

For a composite system of Alice ($a$) and Bob ($b$), the wave function is $\psi(a,b)$.

\begin{definition}[Product state]
The wave function factorizes:
\[
    \psi(a,b) = \psi_A(a)\,\phi_B(b).
\]
\end{definition}

An \textbf{entangled state} is one where no such factorization exists.

\subsection{3.1 Observables on composite systems}
Alice's observable $L$ has matrix elements
\[
    L_{a'b',\,ab} = L_{a'a}\,\delta_{b'b} \qquad\text{(identity on Bob)}.
\]
Bob's observable $M$:
\[
    M_{a'b',\,ab} = \delta_{a'a}\,M_{b'b} \qquad\text{(identity on Alice)}.
\]

% ============================================================
\section{4. The Density Matrix}
% ============================================================

When Alice has access only to her subsystem and cannot perform experiments on Bob, she needs a reduced description.

\begin{definition}[Alice's density matrix]
\[
    \boxed{\rho^{(A)}_{a,a'} = \sum_b \psi^*(a',b)\,\psi(a,b).}
\]
Bob's variables have been \textbf{summed over} (``traced out'').
\end{definition}

\subsection{4.1 Properties}
\begin{itemize}[nosep]
    \item $\rho^{(A)}$ is \textbf{Hermitian}: $(\rho^{(A)}_{a,a'})^* = \rho^{(A)}_{a',a}$.
    \item $\operatorname{Tr}\rho^{(A)}=1$ (normalization).
    \item For any Alice observable $L$: $\expval{L} = \operatorname{Tr}(L\,\rho^{(A)}) = \sum_{a,a'} L_{a'a}\,\rho^{(A)}_{a,a'}$.
    \item $\rho^{(A)}$ contains \emph{all} statistical information Alice can ever extract from her subsystem alone.
\end{itemize}

\subsection{4.2 Product state}
If $\psi(a,b)=\psi_A(a)\phi_B(b)$, then
\[
    \rho^{(A)}_{a,a'} = \psi_A(a)\,\psi_A^*(a') \cdot \underbrace{\sum_b |\phi_B(b)|^2}_{=1} = \psi_A(a)\,\psi_A^*(a').
\]
This is a \textbf{pure-state} density matrix.  It has one eigenvalue $=1$ (eigenvector $\psi_A$) and all others $=0$.

\noindent\textbf{Proof:} Act with $\rho^{(A)}$ on a vector $\chi(a)$:
\[
    \sum_{a'} \rho^{(A)}_{a,a'}\chi(a') = \psi_A(a)\sum_{a'}\psi_A^*(a')\chi(a') = \psi_A(a)\cdot\braket{\psi_A}{\chi}.
\]
If $\chi=\psi_A$: result is $\psi_A(a)\cdot 1 = \psi_A(a)$ (eigenvalue $+1$).  If $\chi\perp\psi_A$: result is $0$ (eigenvalue $0$).  So $(\rho^{(A)})^2 = \rho^{(A)}$---a \textbf{projection operator}.

\subsection{4.3 Entangled state}
For the singlet $\ket{\Psi}=\frac{1}{\sqrt{2}}(\ket{\uparrow\downarrow}-\ket{\downarrow\uparrow})$:
\[
    \rho^{(A)} = \frac{1}{2}\begin{pmatrix}1&0\\0&1\end{pmatrix} = \frac{1}{2}\mathbf{1}.
\]
This is a \textbf{maximally mixed} state: both eigenvalues are $\frac{1}{2}$.  Alice's subsystem appears completely random---she has \emph{zero} information about her spin alone, even though the composite state is a perfectly known pure state.

\subsection{4.4 Entanglement criterion via the density matrix}
The eigenvalues of $\rho^{(A)}$ quantify entanglement:
\begin{itemize}[nosep]
    \item \textbf{Product state (no entanglement):} one eigenvalue $=1$, rest $=0$.  $\operatorname{Tr}(\rho^{(A)})^2 = 1$.
    \item \textbf{Maximal entanglement:} all eigenvalues equal ($=1/n$ for dimension $n$).  $\operatorname{Tr}(\rho^{(A)})^2 = 1/n$.
    \item \textbf{Partial entanglement:} intermediate case.  $1/n < \operatorname{Tr}(\rho^{(A)})^2 < 1$.
\end{itemize}

% ============================================================
\section{5. Correlations Revisited}
% ============================================================

\begin{definition}[Correlation]
\[
    C(L,M) = \expval{LM} - \expval{L}\expval{M}.
\]
\end{definition}

For the singlet:
\begin{align}
    \expval{\sigma_z} &= 0, & \expval{\tau_z} &= 0, & \expval{\sigma_z\tau_z} &= -1, \\
    \expval{\sigma_x} &= 0, & \expval{\tau_x} &= 0, & \expval{\sigma_x\tau_x} &= -1.
\end{align}
Every pair of matching components has correlation $-1$.  This is \textbf{maximal entanglement}: knowing one spin's measurement outcome determines the other's with certainty (opposite result).

\subsection{5.1 Criterion for entanglement}
A composite state is entangled if and only if there exists \emph{at least one} pair of subsystem observables $(L,M)$ with $C(L,M)\neq 0$.  Equivalently, the reduced density matrix $\rho^{(A)}$ is \emph{not} a pure state (i.e.\ $\operatorname{Tr}(\rho^{(A)})^2 < 1$).

% ============================================================
\section{6. Measurement as Entanglement}
% ============================================================

A measurement is the process of entangling a system with an apparatus.  Model: the apparatus has two states, $\ket{0}$ (neutral/ready) and $\ket{1}$ (detected).  The interaction rule:
\[
    \ket{\uparrow}\ket{0} \to \ket{\uparrow}\ket{1}, \qquad
    \ket{\downarrow}\ket{0} \to \ket{\downarrow}\ket{0}.
\]
For a general spin state $\alpha\ket{\uparrow}+\beta\ket{\downarrow}$, the combined initial state $\bigl(\alpha\ket{\uparrow}+\beta\ket{\downarrow}\bigr)\ket{0}$ evolves to:
\[
    \alpha\ket{\uparrow}\ket{1} + \beta\ket{\downarrow}\ket{0}.
\]
This is an \textbf{entangled state} between spin and apparatus.  By looking at the apparatus, you learn about the spin---that is exactly what a measurement is.

If you include the apparatus in the quantum description, there is no ``collapse''---only the development of entanglement via unitary evolution.  Collapse is what you say when you choose \emph{not} to include the apparatus in the quantum system.  This hierarchy (system $\to$ apparatus $\to$ observer $\to$ \ldots) is \textbf{Wigner's friend} scenario---you can draw the line anywhere, and the results are consistent.

% ============================================================
\section{7. Locality and Entanglement}
% ============================================================

Suppose Alice and Bob share an entangled state and separate to distant locations.

\noindent\textbf{Key result:} Any operation Bob performs on his subsystem alone (measurement, unitary transformation, etc.) has \emph{no effect} on Alice's density matrix $\rho^{(A)}$.  Therefore:
\begin{itemize}[nosep]
    \item Bob cannot send a message to Alice by manipulating his spin.
    \item All of Alice's measurement statistics remain unchanged.
    \item This is the \textbf{no-signaling} property of quantum mechanics.
\end{itemize}

\noindent The correlations in the singlet state are analogous to the penny-and-dime story: if Alice sees heads, she instantly knows Bob has tails.  But this is just a pre-established correlation, not a signal.  The genuinely quantum feature (Bell's theorem, next lecture) is that these correlations cannot be reproduced by any local classical model.

% ============================================================
\section{8. Key Conceptual Points}
% ============================================================

\begin{enumerate}[nosep]
    \item A composite system can be in a definite pure state $\ket{\Psi}$ while each subsystem has \emph{no} definite state of its own---only a density matrix.  You can ``know everything about the whole and nothing about the parts.''
    \item This was Einstein, Podolsky, and Rosen's (EPR) central puzzle.
    \item The density matrix formalism makes it precise: the subsystem's statistical description is obtained by tracing over the other subsystem's degrees of freedom.
    \item Entanglement is \emph{not} the same as ``sharing a state''---it is the failure of the wave function to factorize.
    \item Measurement is the development of entanglement between system and apparatus---collapse is the description from the perspective of the system alone.
\end{enumerate}

% ============================================================
\section{9. Summary}
% ============================================================

\begin{enumerate}[nosep]
    \item The \textbf{wave function} $\psi(a,b,\dots)=\braket{a,b,\dots}{\Psi}$ is the set of expansion coefficients in a chosen basis.
    \item The \textbf{density matrix} $\rho^{(A)}_{a,a'}=\sum_b \psi^*(a',b)\psi(a,b)$ describes a subsystem by tracing over inaccessible degrees of freedom.
    \item In a product state, $\rho^{(A)}$ is a pure-state projector (one eigenvalue $=1$); in an entangled state, it is mixed (multiple nonzero eigenvalues).
    \item \textbf{Maximal entanglement}: all eigenvalues of $\rho^{(A)}$ are equal ($=1/n$).
    \item Nonzero correlations $C(L,M)\neq 0$ between subsystem observables signal entanglement.
    \item \textbf{Measurement is entanglement}: the apparatus becomes correlated with the system via unitary evolution.  Collapse is the description when the apparatus is not included.
    \item \textbf{Locality}: Bob's operations cannot change Alice's density matrix---no signaling.
\end{enumerate}

\end{document}
